\hypertarget{glossary}{}
\hypertarget{terminology-glossary}{%
\chapter{Terminology glossary}\label{terminology-glossary}}

\hypertarget{glossary-grid}{}
\textbf{Alpha}\\
Regression intercept: average return unexplained by the included linear
factors.

\textbf{Beta}\\
Factor loading measuring conditional sensitivity to a factor.

\textbf{Excess return}\\
Portfolio return minus the risk-free return.

\textbf{FF3}\\
Fama-French market, size and value factor model.

\textbf{FF5}\\
FF3 augmented with profitability and investment factors.

\textbf{Mkt-RF}\\
Market portfolio return minus the risk-free rate.

\textbf{SMB}\\
Small-minus-big factor associated with firm size.

\textbf{HML}\\
High-minus-low book-to-market value factor.

\textbf{RMW}\\
Robust-minus-weak profitability factor.

\textbf{CMA}\\
Conservative-minus-aggressive investment factor.

\textbf{Basis point}\\
One hundredth of one percentage point; 100 bps = 1\%.

\textbf{Balanced panel}\\
Every portfolio has an observation at every included date.

\textbf{Analytical grain}\\
The entity-time key represented by one row; here, one portfolio-month.

\textbf{Provenance}\\
Documented origin, date, version, transformation and identity of data.

\textbf{SHA-256}\\
Cryptographic digest used to identify exact file bytes.

\textbf{OLS}\\
Ordinary least squares regression.

\textbf{Self-financing portfolio}\\
A long-short portfolio whose long position is financed by its short
position, so net initial investment is zero in the idealised
construction.

\textbf{Residual}\\
Observed outcome minus fitted model outcome.

\textbf{Influence weight}\\
Coefficient showing how one observation\textquotesingle s disturbance
contributes linearly to an estimator such as alpha.

\textbf{Heteroskedasticity}\\
Non-constant conditional residual variance.

\textbf{Autocorrelation}\\
Dependence between observations at different times.

\textbf{Cross-sectional dependence}\\
Same-date dependence among different portfolios.

\textbf{HAC}\\
Heteroskedasticity- and autocorrelation-consistent covariance
estimation.

\textbf{Newey-West}\\
A kernel-weighted HAC covariance estimator.

\textbf{Long-run covariance}\\
Sum of contemporaneous and lagged covariance matrices governing the
asymptotic variance of a dependent sample mean.

\textbf{Bartlett kernel}\\
Linearly declining weights applied to lagged covariance terms.

\textbf{Lag}\\
Time separation between observations in dependence calculations.

\textbf{Positive semi-definite}\\
A covariance property ensuring non-negative variance in every direction.

\textbf{Condition number}\\
Numerical sensitivity of a matrix solve to perturbations.

\textbf{Wald test}\\
Quadratic-form test of parameter restrictions using an estimated
covariance.

\textbf{GRS test}\\
Classical joint test that all factor-model alphas are zero under IID
Gaussian assumptions.

\textbf{Null hypothesis}\\
The restriction evaluated by a statistical test.

\textbf{p-value}\\
Tail probability of evidence at least as extreme under the null model.

\textbf{False discovery rate}\\
Expected proportion of false rejections among all rejections.

\textbf{BH procedure}\\
Benjamini-Hochberg step-up multiple-testing correction.

\textbf{q-value}\\
Multiplicity-adjusted evidence associated with FDR control.

\textbf{Empirical Bayes}\\
Bayesian-form shrinkage with prior hyperparameters estimated from the
observed ensemble.

\textbf{Prior}\\
Distribution describing the cross-sectional alpha population before each
noisy estimate is combined with it.

\textbf{Posterior}\\
Updated distribution after combining prior structure with data.

\textbf{Hyperparameter}\\
Parameter of the prior distribution, such as \ensuremath{\mu } or \ensuremath{\tau }.

\textbf{Marginal likelihood}\\
Likelihood after integrating out latent parameters; used here to
estimate \ensuremath{\mu } and \ensuremath{\tau }.

\textbf{Profile likelihood}\\
Likelihood in one parameter after replacing nuisance parameters by their
conditional maximisers.

\textbf{Partial pooling}\\
Shrinking unit estimates toward a shared centre without forcing them to
be equal.

\textbf{Winner\textquotesingle s curse}\\
Upward noise bias in an estimate selected because it is the largest.

\textbf{Block bootstrap}\\
Resampling contiguous observations to preserve local time dependence.

\textbf{Circular block}\\
A block allowed to wrap from the sample end to the beginning.

\textbf{Common-date resampling}\\
Applying identical resampled dates to every portfolio.

\textbf{Rank interval}\\
Empirical uncertainty range for an item\textquotesingle s rank.

\textbf{Rolling window}\\
Fixed-length historical sample that moves through time.

\textbf{Expanding window}\\
Historical sample that retains all available data up to each cutoff.

\textbf{Look-ahead}\\
Use of information unavailable at the decision time.

\textbf{Leakage}\\
Any path by which evaluation information influences training,
preprocessing or selection.

\textbf{Frozen beta}\\
A factor loading fixed at the training cutoff during future evaluation.

\textbf{Spearman correlation}\\
Correlation of ranks; a measure of monotone ordering agreement.

\textbf{Fisher transformation}\\
atanh transformation used when aggregating correlations.

\textbf{Quintile spread}\\
Mean outcome of the top fifth minus mean outcome of the bottom fifth.

\textbf{Oracle}\\
Look-ahead benchmark used as an upper bound, not an implementable
strategy.

\textbf{Jaccard index}\\
Set overlap measured as intersection divided by union.

\textbf{Persistence}\\
Stability of scores or ranks across time.

\textbf{Mobility}\\
Movement between rank or quintile states over time.

\textbf{Jarque-Bera}\\
Residual normality diagnostic based on skewness and kurtosis.

\textbf{Ljung-Box}\\
Portmanteau test for serial autocorrelation.

\textbf{ARCH-LM}\\
Test for autoregressive conditional heteroskedasticity.

\textbf{SFA}\\
Stochastic frontier analysis separating symmetric noise and one-sided
shortfall.

\textbf{Composite error}\\
Disturbance formed from multiple latent components; in production SFA,
\ensuremath{\varepsilon }=v\ensuremath{-}u.

\textbf{Half-normal}\\
Distribution of the absolute value of a zero-mean normal variable.

\textbf{Truncated normal}\\
A normal distribution conditioned to lie inside a specified interval.

\textbf{Inverse Mills ratio}\\
Ratio \ensuremath{\phi }(z)/\ensuremath{\Phi }(z), appearing in moments of a one-sided truncated normal.

\textbf{Likelihood}\\
Observed-data density viewed as a function of unknown parameters.

\textbf{Maximum likelihood}\\
Parameter estimation by maximising the sample likelihood or
log-likelihood.

\textbf{Score}\\
Gradient of the log-likelihood with respect to its parameters; distinct
here from a portfolio performance score.

\textbf{Mixture distribution}\\
Probability distribution formed by weighting component distributions;
the SFA boundary reference mixes \ensuremath{\chi }\textsuperscript{2}\textsubscript{0} and
\ensuremath{\chi }\textsuperscript{2}\textsubscript{1}.

\textbf{Identification}\\
Property that distinct parameter values imply distinct observable
probability distributions.

\textbf{Boundary test}\\
Test where the null parameter lies on the edge of its permitted space.

\textbf{Model risk}\\
Conclusion sensitivity to benchmark, assumptions, specification or
tuning.

\textbf{Manifest}\\
Machine-readable record of inputs, settings, environment and outputs.

\textbf{NOT COMPARABLE}\\
Status indicating that exact equality is not a valid check for the
available versions.

